\documentclass[10pt,a4paper]{amsart}
\usepackage{amsmath,amssymb,amsthm,graphicx}
\usepackage{amsmath}
\usepackage{amssymb}
\usepackage{hyperref}
\usepackage{amsmath,amssymb,amsthm,enumerate,mathabx}

 

\newtheorem{Theorem}{Theorem}[section]
\newtheorem{Proposition}[Theorem]{Proposition}
\newtheorem{Lemma}[Theorem]{Lemma}
\newtheorem{Corollary}[Theorem]{Corollary}
\newtheorem{Theorem1}{Theorem}
 
\theoremstyle{definition}
\newtheorem{Definition}[Theorem]{Definition}
\newtheorem{Fact}[Theorem]{Fact}
\newtheorem{Conjecture}[Theorem]{Conjecture}
\newtheorem{Example}[Theorem]{Example}
\newtheorem{Remark}[Theorem]{Remark}
\newtheorem{Question}{Question}
\newtheorem{Claim}{Claim}

\DeclareMathOperator{\cl}{cl_S}
\DeclareMathOperator{\dcl}{dcl}
\DeclareMathOperator{\tp}{tp}
\DeclareMathOperator{\Aut}{Aut}
\DeclareMathOperator{\conv}{conv} 
\DeclareMathOperator{\dom}{dom}
\DeclareMathOperator{\acc}{acc}

\usepackage[margin=23mm]{geometry}

\title{Vaught's conjecture for theories of discretely ordered structures}
\author{Predrag Tanovi\'c}
\thanks{This research was supported by the Science Fund of the Republic of 
Serbia, Grant No. 7750027: Set-theoretic, model-theoretic and 
Ramsey-theoretic phenomena in mathematical structures: similarity and 
diversity–SMART}
\address{Mathematical Institute SANU, Belgrade, Serbia} 
\email{tane [at] mi.sanu.ac.rs}
\usepackage{cleveref}
\crefname{originallemma}{original lemma}{original lemmas}
\Crefname{originallemma}{Original lemma}{Original lemmas}
\makeatletter
\newlabel{Lemma_basic_simple}{{1.4}{source}{}{sourcepointer-1}{}}
\newlabel{Lemma_basic_simple@cref}{{[originallemma][0][]1.4}{[1][0][]source}{}{}{}}
\makeatother
\title{A relatively definable function on the realisations of a simple type is a power of the successor}
\author{Predrag Tanovic}
\date{arxiv:2212.13605; CC BY 4.0}
\begin{document}
\maketitle
\noindent\textit{Editorial note.} All delivered mathematical text is the author's own, copied from the pinned publisher source sources/vaught-discrete-orders-2212.13605v3/source0.tex; no statement or proof step was rewritten, reordered or removed.\par\medskip
\noindent\textit{Editorial note.} Preamble provenance: the retained preamble is the source's own layout adapter, if any, plus the author's macro block at lines 4-40 of sources/vaught-discrete-orders-2212.13605v3/source0.tex (everything between the publisher class metadata and \textbackslash{}begin\{document\}, comments excluded). Nothing was added, no macro was invented, and the publisher document class is deliberately not reproduced; the excerpt serves only as preamble, never as a span coordinate.\par\medskip
\noindent\textit{Editorial note.} Metadata provenance: source\_id and source\_version were taken from SOURCE\_RECORD.json; the title from SOURCE\_RECORD.json; the author list from the licence evidence (its authors entry); the source URL from the first non-empty of source\_url, source\_urls[0] and the evidence's fixed\_version\_url/license\_url. The licence evidence file used by this card is sources/vaught-discrete-orders-2212.13605v3/LICENSE\_EVIDENCE.json, recorded in license.evidence\_path so that the binding target is explicit.\par\medskip
\noindent\textit{Editorial note.} External references are navigation pointers only; every deliberately kept step of the author's argument stays inside the counted spans. The displayed numbers of those pointers are the author's own numbering under the author's \textbackslash{}newtheorem declarations (the source has no subsections, so numbering is section.index), computed mechanically by the lane's scanner.\par\medskip
\noindent\textit{Editorial note.} The single reference inside the span is Lemma\_basic\_simple, recorded as a navigation stub of kind lemma; the step that consumes it stays inside the counted span.\par\medskip
\noindent\textit{Editorial note.} The counted spans contain no bibliography citation, so no bibliography excerpt is attached; the card points at the automatically generated per-card licence reuse record because the staged licence evidence carries no status field.\par\medskip
\noindent\textit{Allowed standard prerequisites (editorial).}\begin{itemize}
\item The clauses of the basic lemma on simple types, in particular that every relatively definable nonempty subset of p(M) has a minimum. Same source, Lemma labelled Lemma\_basic\_simple, clause (c); invoked inside the counted proof at lines 224 and 229. Supplies the minimum used for injectivity and the contradiction with non-minimality used for order preservation.
\item Aleph\_0-saturation of the ambient model and compactness. Used inside the counted proof at lines 223, 231 and 236. Gives surjectivity of the extension of f and lets the interval-wise behaviour found on a tail of C be transferred to the whole type p.
\end{itemize}\medskip
\ifdefined\hypertarget\hypertarget{sourcepointer-1}{}\fi
\noindent\textit{Source reference.} Original lemma 1.4: sources/vaught-discrete-orders-2212.13605v3/source0.tex, source lines 224--224. Navigation pointer for \textbackslash{}ref\{Lemma\_basic\_simple\} used inside the counted span (line 224); the argument that consumes it is retained in the span.\par\medskip
\begin{Lemma}\label{Lemma_f_is_Sn}
Suppose that  $p\in S_1(T)$ is a simple type, as witnessed by $(C,<)$ and $(D,<)$.  If $f:p(M)\to p(M)$ is a relatively definable function, then $f=S^k$ for some $k\in\mathbb Z$, where $S$ is the successor function on $(p(M),<)$.
\end{Lemma}

\begin{proof}Let $c_0=\min C$, $c_n=S^n(c_0)$ for $n\in\mathbb N$, and  suppose that $(D,<)$ is discrete.   
First, we will prove that $f$ is an automorphism of $(p(M),<)$. Choose a definable function $\hat f:D_1\to D_2$ extending $f$. Since $p(x)\cup\{\hat f(x)=a\}$ is consistent for some (equivalently all) $a\in p(M)$ and $M$ is $\aleph_0$-saturated, $f$ is surjective.
To verify the injectivity of $f$, let us consider the set $\hat f^{-1}(\{f(a)\})\cap [c_0,a]\subset p(M)$. It is an $a$-definable subset of $p(M)$ so, by Lemma \ref{Lemma_basic_simple}(c), has minimum, say $a'$; clearly,   $a'=\min(\hat f^{-1}(\{f(a')\})\cap [c_0,a'])$. Since  $a'\in p(M)$, the formula  $x=\min(\hat f^{-1}(\{f(x)\})\cap [c_0,x])$ belongs to $p(x)$; it follows that $f$ is injective. 
In particular, the inverse function $f^{-1}$ is relatively definable, so after possibly shrinking $D_1$ and $D_2$, we will assume that $\hat f$ is a bijection.
It remains to show that $f$ (equivalently $f^{-1}$) is order-preserving. 
Note that at least one of the formulae $\hat f(x)\leq x$ and $\hat f^{-1}(x)\leq x$ belongs to $p(x)$. We will from now on assume $(\hat f(x)\leq x)\in p(x)$, the proof in the other case is similar.
Suppose, by way of contradiction, that $f$ is not order-preserving. Then the set $D(a)$ defined by $\hat f(a)\leq \hat f(x)<x<a$ is nonempty (for all $a\in p(M)$). Observe that  $D(a)\subset [f(a),a]\subset p(M)$,  so  $\emptyset\neq D(b)\subset p(M)$ holds for all $b\in D(a)$. Choose $b'\in D(b)$. Then:  
$$\hat f(a)\leq \hat f(b)\leq \hat f(b')<b'<b<a$$ witnesses non-minimality of $b$ in $D(a)$. Hence no element of $D(a)$ is minimal, contradicting Lemma \ref{Lemma_basic_simple}(c). Therefore, $f$ is an automorphism of $(p(M),<)$. 

By compactness, we can modify $\hat f$ so that it is strictly increasing. In fact, in what follows we will view $\hat f$ as a definable, increasing, partial function  $\hat f\subset D\times D$. Since $\hat f$ is increasing and $\hat f(p(M))=p(M)$, $\hat f(C)\subseteq C$.


Let $\phi(x,y)$ be a formula implying $x<y$ and expressing  ``$[x,y]\subset \dom (\hat f)$ and $\hat f\restriction[x,y]$ is an order-isomorphism of $(D,<)$-intervals $[x,y]$ and $[\hat f(x),\hat f(y)]$''. Then:
$$ p(x)\cup p(y)\cup\{x<y\}\vdash \phi(x,y).$$
By compactness, there is a formula $\theta(x)\in p(x)$   such that:
\begin{equation}
\models  \forall x\forall y  (\theta(x)\land \theta(y)\land x<y\rightarrow \phi(x,y)).
\end{equation}
Since $p$ is a simple type and $\theta(x)\in p$, the set $\lnot\theta(C)$ is finite, so there is a $n_0\in \omega$ such that $\{c_i\mid i\geq n_0\}\subset \theta(C)$. For all $n\geq n_0$, by (1),  $f$ maps the (two-element) interval $[c_n,c_{n+1}]$ onto the interval $[\hat f(c_n),\hat f(c_{n+1})]$, so  
$\hat f(c_n)$ and $\hat f(c_{n+1})$ are consecutive elements of $D$, i.e $\hat f(S(c_n))=S(\hat f(c_n))$ (where $S$ is the succesor function on $(D,<)$). By induction we easily get  $\hat f(S^m(c_{n_0}))=S^m(\hat f(c_{n_0}))$ for all $m\in \omega$. Since $\hat f(c_{n_0})\in C$, there is a $k\in \mathbb Z$ with $\hat f(c_{n_0})=S^k(c_{n_0})$. Then for all $n\geq n_0$: 
$$\hat f(c_n)=\hat f(S^{n-n_0}(c_{n_0}))=    S^{n-n_0}(\hat f(c_{n_0}))=S^{n-n_0+k}(c_{n_0})=S^k(S^{n-n_0}(c_{n_0}))=S^k(c_n).$$ 
Hence  $(\hat f(x)=S^k(x))\in \tp(c_n)$ for  all $n\geq n_0$.  Since $p$ is a $C$-type, $(\hat f(x)=S^k(x))\in p(x)$.
\end{proof}

\end{document}
